% Clase 25 --- El vector de Poynting y la caja de largo c*Delta t (§1.3 y §3).
% (a) la terna E, B, S: la direccion sale sola del producto vectorial.
% (b) el paralelepipedo del argumento: TODA la energia que esta adentro cruza la
%     cara A en Delta t, porque en ese tiempo la onda avanza justo ese largo.
\begin{center}
\begin{tikzpicture}[scale=0.92, transform shape]

% ---------------- (a) la terna ----------------
\begin{scope}
  \draw[figvec] (0,0) -- (0,2.1);
  \node[figlbl, figblue, anchor=east] at (-0.08,1.9) {$\mathbf{E}$ ($\hat y$)};
  \draw[figvecres] (0,0) -- (-1.25,-0.88);
  \node[figlbl, figred, anchor=north east] at (-1.1,-0.82) {$\mathbf{B}$ ($\hat z$)};
  \draw[figvec, figamber, line width=1.3pt] (0,0) -- (2.3,0);
  \node[figlbl, figamber, anchor=north] at (1.5,-0.12) {$\mathbf{S}$ ($\hat x$)};

  \node[figlbl, anchor=north, align=center] at (0.4,-1.75)
    {\textbf{(a)} $\mathbf{S} = \dfrac{1}{\mu_0}\mathbf{E}\times\mathbf{B}$:
     \; $\hat y \times \hat z = \hat x$};
\end{scope}

% ---------------- (b) el paralelepípedo ----------------
\begin{scope}[xshift=6.3cm]
  % cara de entrada (área A)
  \fill[figblue!10] (0,-0.9) -- (0,1.3) -- (0.85,1.9) -- (0.85,-0.3) -- cycle;
  \draw[figline] (0,-0.9) -- (0,1.3) -- (0.85,1.9) -- (0.85,-0.3) -- cycle;
  \node[figlbl, figblue, anchor=east] at (-0.12,0.2) {$A$};

  % aristas hacia el fondo
  \draw[figline] (0,1.3) -- (4.2,1.3);
  \draw[figline] (0.85,1.9) -- (5.05,1.9);
  \draw[figline] (0,-0.9) -- (4.2,-0.9);
  \draw[figaux] (0.85,-0.3) -- (5.05,-0.3);
  \draw[figline] (4.2,-0.9) -- (4.2,1.3) -- (5.05,1.9) -- (5.05,-0.3);
  \draw[figaux] (4.2,-0.9) -- (5.05,-0.3);

  % la onda avanza
  \draw[figvec, figamber, line width=1.3pt] (1.2,0.5) -- (3.0,0.5);
  \node[figlbl, figamber, anchor=south] at (2.1,0.58) {$\mathbf{S}$};

  \draw[figgray, line width=0.6pt,
        {Latex[length=1.8mm,width=1.3mm]}-{Latex[length=1.8mm,width=1.3mm]}]
    (0,-1.25) -- (4.2,-1.25);
  \node[figlblaux, anchor=north] at (2.1,-1.3) {$c\,\Delta t$};

  \node[figlbl, anchor=north, align=center] at (2.1,-1.75)
    {\textbf{(b)} $\Delta U = u\,(A\,c\,\Delta t)
     \;\Rightarrow\; \dfrac{P}{A} = u\,c$};
\end{scope}

\end{tikzpicture}\\[0.5em]
{\small\itshape La dirección de $\mathbf{S}$ no hay que postularla: sale sola de
la definición. Al ser un producto vectorial es perpendicular a $\mathbf{E}$ y a
$\mathbf{B}$, y como en una onda plana esos dos generan el plano transversal, lo
único que queda es la dirección de propagación. El módulo se justifica con la
caja de (b): si se la hace de largo $c\,\Delta t$, en el tiempo $\Delta t$ la
onda avanza exactamente ese largo, de modo que \textbf{toda} la energía que
había adentro terminó cruzando la cara $A$ ---ni más, porque nada entró desde
más atrás, ni menos. Dividiendo por $\Delta t$ y por $A$ queda $P/A = u\,c$, que
es justo lo que vale $|\mathbf{S}| = EB/\mu_0 = cB^2/\mu_0 = c\,u$, usando
$E = cB$ y que las dos densidades de energía son iguales. El $\Delta t$ hay que
tomarlo \textbf{mucho mayor que el período}: la onda oscila, la energía que
cruza no es pareja instante a instante, y lo que se obtiene es un promedio ---que
es además lo único que cualquier instrumento puede medir cuando la luz oscila
$10^{15}$ veces por segundo.}
\end{center}
